\section{Alternative Hessian proofs}
\label{sec:alternative}

The main proof obtains scalar concavity and matrix log-determinant
convexity from the same contraction argument. Here we give direct
Hessian proofs of the two analytic statements. They provide an
alternative route to \Cref{lem:scalar,lem:logdet}; the entropy integral,
Schur complement, sign averaging, and rectangular embedding in
\Cref{sec:block} are unchanged.

\subsection{The Hessian of the scalar deficit}

Recall the function $F$ from \eqref{eq:block-F}. Its homogeneity follows
by the scaling cancellation already proved in \Cref{lem:scalar}.
We verify concavity directly on the probability simplex, where
$F(x,y,z)=H(x,y,z)-g(yz)$, using the derivatives
\eqref{eq:g-derivatives}. For $0<u<1/4$, write
$t=\sqrt{1-4u}$ and $L=\atanh(t)/t$.

\begin{proof}[Alternative proof of \Cref{lem:scalar}]
In the simplex interior, set $x=1-y-z$ and $u=yz$.
The negative Hessian of $F(1-y-z,y,z)$ is
\[
 J=\begin{pmatrix}
 \dfrac1y+\dfrac1x+z^2g''(u)&\dfrac1x+g'(u)+ug''(u)\\[5pt]
 \dfrac1x+g'(u)+ug''(u)&\dfrac1z+\dfrac1x+y^2g''(u)
 \end{pmatrix}.
\]
Expanding its determinant and substituting \eqref{eq:g-derivatives}
and $x+y+z=1$ yields
\begin{align}
 \det J
 &=\frac1{xu}+\frac{y+z-4u}{x}g''(u)
   -\frac2xg'(u)-g'(u)^2-2ug'(u)g''(u)\notag\\
 &=\frac{1-4uL^2}{ut^2}>0.
 \label{eq:block-scalar-determinant}
\end{align}
Indeed, for $0<t<1$,
\[
 \atanh t<\frac{t}{\sqrt{1-t^2}},
\]
as both sides vanish at zero and their derivatives are respectively
$(1-t^2)^{-1}$ and $(1-t^2)^{-3/2}$. Since $4u=1-t^2$, this proves
the strict inequality in \eqref{eq:block-scalar-determinant}.
Furthermore,
\[
 \begin{pmatrix}y&z\end{pmatrix}J\begin{pmatrix}y\\z\end{pmatrix}
 =\frac{y+z}{x}
   +4u\left(ug''(u)+\frac12g'(u)\right)>0.
\]
A real symmetric $2\times2$ matrix with positive determinant has
eigenvalues of the same sign; this positive quadratic form forces
both to be positive. Hence $F$ is concave in the simplex interior,
and continuity extends concavity to the closed simplex.

To prove \eqref{eq:block-superadditive}, omit zero triples and put
$m_i=x_i+y_i+z_i$, $m=\sum_i m_i$. For $m>0$, homogeneity and
simplex concavity give
\[
 \sum_i F(x_i,y_i,z_i)
 =m\sum_i\frac{m_i}{m}
      F\!\left(\frac{x_i}{m_i},\frac{y_i}{m_i},\frac{z_i}{m_i}\right)
 \le F\!\left(\sum_i x_i,\sum_i y_i,\sum_i z_i\right).
\]
If $m=0$, all terms vanish.

\end{proof}

\subsection{The Hessian of the log-determinant difference}

The contraction proof avoids choosing singular-vector coordinates.
The following calculation instead diagonalizes the fixed matrix and
expresses the second derivative as a sum of nonnegative squares.
It proves \Cref{lem:logdet} directly for the same domain, including
singular $Z$.

\begin{proof}[Alternative proof of \Cref{lem:logdet}]
At any point $(D_0,E_0)$, use the fixed congruences
\[
 D=D_0^{1/2}\widetilde D D_0^{1/2},\qquad
 E=E_0^{1/2}\widetilde E E_0^{1/2},\qquad
 W=D_0^{-1/2}ZE_0^{-1/2}.
\]
The determinant factors cancel to give
$G_Z(D,E)=G_W(\widetilde D,\widetilde E)$.
These invertible linear changes preserve affine lines. Independent
unitary changes in the two spaces then put $W$ in singular-value form
$Y=\diag(v_1,\ldots,v_N)$ with $v_i\ge0$. It therefore suffices to
prove nonnegativity of the second derivative at $(I,I)$ for this $Y$
and arbitrary Hermitian directions $P,Q$.

Set $R=(I+YY^*)^{-1}$ and $N_0=P-YQY^*$.
Using $(I+sQ)^{-1}=I-sQ+s^2Q^2+O(s^3)$ gives
\begin{equation}\label{eq:block-G-hessian}
 \left.\frac{\mathrm d^2}{\mathrm ds^2}
 G_Y(I+sP,I+sQ)\right|_{s=0}
 =\Tr\bigl(P^2+2RYQ^2Y^*-RN_0RN_0\bigr).
\end{equation}
For completeness, the log-determinant expansion used here is
\[
 \ln\det(H+sK+s^2L_0)
 =\ln\det H+s\Tr(H^{-1}K)
 +s^2\left(\Tr(H^{-1}L_0)
       -\frac12\Tr(H^{-1}KH^{-1}K)\right)+O(s^3),
\]
for $H>0$ and Hermitian $K,L_0$. It follows by factoring $H^{1/2}$
and applying the scalar logarithm expansion to the eigenvalues of the
resulting Hermitian increment. No commutativity is assumed.

Write $r_i=(1+v_i^2)^{-1}$ and $c_i=v_i/(1+v_i^2)$.
Expanding the trace in \eqref{eq:block-G-hessian} in entries gives
\begin{align*}
 \sum_{i,j}\bigl[&(1-r_ir_j)|P_{ij}|^2
 +\{2(1-r_i)-(1-r_i)(1-r_j)\}|Q_{ij}|^2\\
 &\hspace{34mm}+2c_ic_j\operatorname{Re}
       (P_{ij}\overline{Q_{ij}})\bigr].
\end{align*}
Since $Q$ is Hermitian, $|Q_{ij}|^2=|Q_{ji}|^2$.
Pairing the $(i,j)$ and $(j,i)$ terms replaces their coefficient by
\[
 \frac12\bigl[2(1-r_i)+2(1-r_j)
              -2(1-r_i)(1-r_j)\bigr]=1-r_ir_j.
\]
This identity also holds when $i=j$. Completing the square therefore
rewrites the second derivative as
\begin{equation}\label{eq:block-G-squares}
 \sum_{i,j}\left[
 (1-r_ir_j-c_ic_j)(|P_{ij}|^2+|Q_{ij}|^2)
       +c_ic_j|P_{ij}+Q_{ij}|^2\right].
\end{equation}
Every coefficient is nonnegative: $c_ic_j\ge0$, and
\[
 1-r_ir_j-c_ic_j
 =\frac{(v_i-v_j)^2+v_iv_j+v_i^2v_j^2}
        {(1+v_i^2)(1+v_j^2)}\ge0.
\]
Thus the second derivative is nonnegative in every Hermitian
direction at every positive definite pair. The domain is convex,
so the asserted joint convexity follows along each line segment.
\end{proof}
